\subsection{How much predicted recovery does an extra place cost?}\label{sec:shortlist-frontier}
A platform need not choose between independent row rankings and an arbitrary coordinated assignment. Given a desired number $r$ of feasible placements, it can give up the smallest possible amount of predicted partner coverage. Let $u_j$ be the number of seats at program $j$, $q_i=\min(k,d_i)$ the number of recommendations for applicant $i$, and $s_{ij}$ a score on a common cardinal scale. When $s_{ij}$ is the probability that $j$ is $i$'s full-list partner, summing the displayed scores gives the predicted number whose partner is shown, without requiring independence across people. Write $D_0$ for rowwise top-$q_i$ and $\tau_i$ for the last retained score. Skip rows with no slots. Write $\nu(D)$ for the maximum number of applicants that the display can assign within the seat capacities.

\begin{theorem}[The price of visible capacity]\label{thm:capacity-frontier}
Suppose all applicants can receive a separately chosen list. For every $0\le r\le\nu(E)$ the minimum loss of total score from $D_0$, over quota-respecting displays $D$ with $\nu(D)\ge r$, is
\[
L(r)=\min_{A:\,|A|=r}\;\sum_{(i,j)\in A}(\tau_i-s_{ij})_+ ,
\]
where $A$ ranges over feasible assignments respecting program capacities. The nonnegative marginal losses $L(r+1)-L(r)$ are nondecreasing. For one recommendation per applicant, $r$ is also a lower bound on the number placed by DA under \emph{every} strict preference and responsive priority order on the display.
\end{theorem}
\begin{proof}
Any admissible display contains an assignment of size $r$. Each of its rows containing an assigned edge $(i,j)$ loses at least $(\tau_i-s_{ij})_+$ relative to its best list. Distinct assigned edges use distinct applicant rows, so these lower bounds add. Conversely, take a minimizing assignment. On each assigned row force its assigned edge onto the list and choose the best $q_i-1$ remaining edges. On every other row take its best list. This attains the sum, including when boundary scores tie.

Send $r$ units of flow from a source through applicants, admissible programs and a sink, with applicant arcs of capacity one and program-to-sink arcs of capacity $u_j$. Only applicant--program arcs carry the displayed replacement costs. Integral flow produces a minimizing assignment. The average of optimal flows of values $r-1$ and $r+1$ is feasible at value $r$. Integrality at that integer value gives $2L(r)\le L(r-1)+L(r+1)$. Deleting one assigned edge shows monotonicity. At quota one, every applicant with an option makes exactly one proposal. Each program holds the smaller of its capacity and its number of applicants, exactly exhausting the display's feasible placement count.
\end{proof}

This uses a standard minimum-cost flow primitive \citep{ahuja1993flows}. Flow-based recommender post-processing previously traded relevance against item exposure counts \citep{antik2017}. Here the constraint is the number of \emph{different people who could simultaneously be assigned}. Showing the same available programs to many applicants can meet an exposure target while leaving too few seats for those particular people. The cost of forcing an edge is its opportunity cost in that person's list. Choosing a maximum-weight assignment on the unadjusted scores optimizes a different objective.

\paragraph{What the platform can guarantee without knowing preferences.}
Under responsive priorities, every stable matching is maximal: an unmatched applicant cannot see a program with a vacant seat. Conversely, if both strict orderings are unknown, every maximal assignment can occur under applicant-proposing DA. To realize one, put each assigned program first on its holder's list and put all designated holders first in its program's priority. Unassigned applicants only see full programs by maximality. Thus the graph alone certifies exactly the size of its smallest maximal assignment in the worst case. In particular it certifies at least $\lceil\nu(D)/2\rceil$: expand each program into its seats. Each edge of a maximum matching either belongs to the maximal matching or meets one of its matched endpoints. Each edge of the maximal matching accounts for at most two edges of the maximum matching. The factor is sharp already at quota two: applicant $a$ sees $x,y$ in that order, $b$ sees only $x$, and $x$ prioritizes $a$. Both can be assigned, but only $a$ is placed by DA. Computing the exact worst-case size is NP-hard for unit-capacity bipartite graphs of maximum degree three by the minimum maximal matching result of \citet{yannakakis1980}. Deciding whether all strict profiles place at least $r$ is coNP-complete, since a smaller maximal assignment certifies failure. These worst-case statements allow preference and priority orderings to vary. For observed orders the realized DA allocation is determined and can be computed directly.
